\subsection{Tipos de Clasificadores para Reconocimiento de Patrones}

El reconocimiento de patrones es una disciplina fundamental en inteligencia
artificial y aprendizaje automático, cuyo objetivo es identificar
regularidades o estructuras en datos para asignarles categorías. Los
clasificadores son algoritmos que implementan esta tarea, y su elección
depende de factores como la naturaleza de los datos, la complejidad del
problema y la necesidad de interpretabilidad. A continuación, se describen
los principales tipos de clasificadores y sus características, con énfasis
en los tres utilizados en esta práctica: $k$-vecinos más cercanos, árboles
de decisión y máquinas de soporte vectorial.

\subsection{Clasificadores No Lineales}

Cuando los datos presentan relaciones complejas, se emplean modelos no
lineales. Los \textbf{árboles de decisión} son un ejemplo: dividen
recursivamente el espacio de características mediante reglas jerárquicas, lo
que los hace altamente interpretables pero propensos al sobreajuste cuando
crecen sin control. Sus principales hiperparámetros (profundidad máxima,
muestras mínimas por nodo, criterio de división) actúan precisamente como
mecanismos de regularización. Otro enfoque es el algoritmo de
\textbf{$k$-vecinos más cercanos (KNN)}, que clasifica instancias según la
mayoría de clases entre sus vecinos más cercanos en el espacio de
características. Aunque es flexible y no requiere entrenamiento explícito
---se trata de un método \textit{lazy} o perezoso---, su costo
computacional crece con el tamaño de los datos y depende fuertemente del
escalado y de la métrica de distancia elegida. Estos métodos son útiles en
reconocimiento de escritura manual, sistemas de recomendación o, como en
esta práctica, en problemas biomédicos con variables heterogéneas.

\subsection{Clasificadores Probabilísticos}

Basados en teoría de probabilidad, estos modelos estiman distribuciones
para predecir clases. El clasificador \textbf{Naive Bayes} es el más
conocido: asume independencia entre características y aplica el teorema de
Bayes para calcular probabilidades posteriores. Aunque simplista, es eficaz
en datos categóricos o textuales, como en análisis de sentimientos o
detección de spam. Las redes Bayesianas, por otro lado, modelan dependencias
entre variables, ofreciendo mayor flexibilidad a costa de complejidad.

\subsection{Máquinas de Soporte Vectorial (SVM)}

Las SVM buscan hiperplanos óptimos para separar clases, incluso en espacios
de alta dimensión. En su versión \textit{lineal}, maximizan el margen entre
las clases y son particularmente efectivas cuando los datos son linealmente
separables o aproximadamente lineales tras un buen escalado. Con
\textit{kernels} (típicamente el gaussiano o RBF) transforman implícitamente
los datos a espacios de mayor dimensión donde la separación no lineal es
posible. Su comportamiento se controla principalmente con dos hiperparámetros:
el parámetro de regularización $C$ ---que pondera el compromiso entre
maximizar el margen y tolerar errores de clasificación--- y, en el caso del
kernel RBF, el parámetro $\gamma$, que define la amplitud de influencia de
cada muestra. Son eficaces en reconocimiento facial o clasificación de
imágenes médicas, pero requieren ajuste cuidadoso de parámetros y
normalización previa de los datos.

\subsection{Redes Neuronales}

Estos modelos, inspirados en el cerebro humano, destacan por su capacidad
para aprender patrones complejos. Los perceptrones multicapa (MLP) son redes
\textit{feed-forward} con capas ocultas no lineales, mientras que las redes
convolucionales (CNN) se especializan en datos \textit{grid-like} como
imágenes. Las redes recurrentes (RNN) manejan secuencias (texto, series
temporales). Aunque son extremadamente poderosas (p.~ej.\ en visión por
computadora o traducción automática), requieren grandes volúmenes de datos,
recursos computacionales y carecen de interpretabilidad directa.

\subsection{La Búsqueda en Cuadrícula como Estrategia de Ajuste}

Cada uno de los clasificadores anteriores expone un conjunto de
hiperparámetros cuya elección puede cambiar significativamente el desempeño.
La búsqueda en cuadrícula (\textit{grid search}) es la estrategia más
sistemática para explorar este espacio: se define una lista de valores
plausibles para cada hiperparámetro y se entrena el modelo para todas las
combinaciones, evaluando cada una mediante validación cruzada. Aunque es
exhaustiva ---y por tanto costosa cuando el número de combinaciones
crece---, garantiza encontrar el óptimo dentro de la cuadrícula definida y
ofrece una visión clara de qué hiperparámetros realmente importan.
